\hmsubsection{Development Methodology}{OAH}{UMA}\label{sec:methodology}

The project was carried out using an iterative and incremental development approach
based on the Scrum framework~\cite{schwaber2020scrum}. The methodology was not applied
in its pure form, but adapted to the size and constraints of the team. This approach
was selected to enable continuous development, frequent feedback, and adaptability
throughout the project lifecycle.

Work was organised into sprints with a duration of one to two weeks. Each sprint
began with planning, where tasks from the backlog were broken into smaller,
actionable items to reduce complexity and enable a try-and-fail cycle. At the end
of each sprint, a retrospective was conducted to identify what worked well, address
challenges, and define improvements for the next iteration. One notable adaptation
was the early decision to abandon MobileNetV2 transfer learning in favour of a
classical computer vision pipeline, following poor real-time performance on the
target hardware.

Task planning and tracking were managed using a digital Trello board, with tasks
categorised into Backlog, In Progress, Testing, and Done. This gave the team a
clear overview of task ownership and progress at all times. Stand-up meetings were
held three times per week on Tuesdays, Wednesdays, and Thursdays, lasting no more
than 15 minutes. Each member reported what they had completed, what they planned to
do next, and whether they had any blockers.

Communication was maintained through a dedicated Microsoft Teams channel, enabling
asynchronous coordination and remote participation when physical attendance was not
possible. Shared project documents, steering documents, and time registration were
stored on a common OneDrive, accessible to all members at all times.

External communication followed a fixed weekly cadence: client meetings every
Monday and supervisor meetings every Thursday. Over the course of the project, the
team participated in two formal presentations, receiving both group and individual
feedback that informed subsequent development decisions.
